\documentclass[a4paper,12pt]{article}
\usepackage[T2A]{fontenc}
\usepackage[utf8]{inputenc}
\usepackage[russian]{babel}
\usepackage{amsmath,amssymb,icomma}
\usepackage[a4paper,left=20mm,right=20mm,top=20mm,bottom=21mm,headheight=25pt,headsep=9pt]{geometry}
\usepackage{microtype,parskip,setspace,enumitem,booktabs,tabularx,fancyhdr,titlesec,xcolor,tikz}
\usetikzlibrary{arrows.meta,positioning,shapes.geometric}
\definecolor{accent}{HTML}{008080}
\definecolor{darkaccent}{HTML}{004D4D}
\definecolor{textgray}{HTML}{333333}
\color{textgray}
\onehalfspacing
\setlength{\parindent}{0pt}
\setlength{\parskip}{5pt}
\setlist[itemize]{leftmargin=1.5em,itemsep=2pt,topsep=3pt}
\setlist[enumerate]{leftmargin=*,itemsep=5pt,topsep=4pt}
\titleformat{\section}{\Large\normalfont\color{darkaccent}}{\thesection.}{0.5em}{}[\vspace{4pt}{\color{accent}\titlerule[0.6pt]}]
\titleformat{\subsection}{\large\normalfont\color{darkaccent}}{}{0pt}{}
\titlespacing*{\section}{0pt}{12pt}{9pt}
\titlespacing*{\subsection}{0pt}{8pt}{3pt}
\pagestyle{fancy}
\fancyhf{}
\fancyhead[L]{\small\setstretch{1} Отношения и пропорции}
\fancyhead[R]{\small\setstretch{1} Ксения Васильченко}
\fancyfoot[C]{\small\thepage}
\renewcommand{\headrulewidth}{0.5pt}
\renewcommand{\headrule}{\hbox to\headwidth{\color{accent}\leaders\hrule height\headrulewidth\hfill}}
\newcommand{\ex}[1]{\par\medskip\textcolor{darkaccent}{#1}\par\smallskip}
\newcommand{\mix}[3]{#1\frac{#2}{#3}}
\begin{document}
\begin{titlepage}
\centering
\vspace*{27mm}
{\Large\color{darkaccent} МАТЕМАТИКА\par}
\vspace{13mm}
{\LARGE Чек-лист\par}
\vspace{10mm}
{\LARGE\color{accent} Отношения и пропорции\par}
{\Large\color{accent} Прямая и обратная пропорциональность\par}
\vspace{14mm}
{\large Школьная математика\par}
\vfill
{\large Ксения Васильченко\par}
\vspace{8mm}
24 ноября 2026 г.\par
\vspace{15mm}
{\small Лёвин Артём Александрович\par эксклюзивно для Ксении Васильченко\par}
\vspace{15mm}
\begin{tikzpicture}
\draw[accent,line width=0.85pt] (0,0) .. controls (2.1,0.8) and (4,-0.8) .. (6,0) .. controls (8,0.8) and (10,-0.8) .. (12,0);
\end{tikzpicture}
\end{titlepage}
\setcounter{page}{1}

\section{Отношение и основное свойство пропорции}
\textbf{Отношение} двух чисел $a$ и $b$ --- это частное $a:b=\dfrac{a}{b}$, где $b\ne0$.

\textbf{Пропорция} --- равенство двух отношений:
\[
 a:b=c:d\quad\Longleftrightarrow\quad \frac{a}{b}=\frac{c}{d},\qquad b\ne0,\ d\ne0.
\]
Числа $a$ и $d$ называют крайними членами, а $b$ и $c$ --- средними.
\begin{center}
\begin{tikzpicture}[every node/.style={font=\normalsize}]
\node (a) at (0,0) {$a$};\node at (0.8,0) {$:$};\node (b) at (1.6,0) {$b$};\node at (2.5,0) {$=$};\node (c) at (3.4,0) {$c$};\node at (4.2,0) {$:$};\node (d) at (5,0) {$d$};
\draw[accent,line width=.7pt] (a.north) to[out=75,in=105] node[above,font=\small] {крайние} (d.north);
\draw[darkaccent,line width=.7pt] (b.south) to[out=-70,in=-110] node[below,font=\small] {средние} (c.south);
\end{tikzpicture}
\end{center}
\textbf{Основное свойство:} произведение крайних членов равно произведению средних. И наоборот: при ненулевых знаменателях равенство произведений означает равенство отношений.
\[
\boxed{\frac{a}{b}=\frac{c}{d}\quad\Longleftrightarrow\quad ad=bc}
\]
\ex{Пример 1. Можно ли составить пропорцию?}
Проверим равенство $3,6:1,2=4,5:1,5$.
\[
3,6\cdot1,5=5,4,\qquad 1,2\cdot4,5=5,4.
\]
Произведения равны, значит, пропорция верна.
\ex{Пример 2. Когда пропорция неверна}
Проверим $0,8:0,2=1,2:0,4$.
\[
0,8\cdot0,4=0,32,\qquad 0,2\cdot1,2=0,24.
\]
Произведения различны, поэтому равенство отношений неверно.

\textcolor{darkaccent}{Проверяйте именно произведения крайних и средних членов, а не отдельные числа.}

\clearpage
\section{Вычисления с десятичными дробями}
При проверке пропорций особенно важна точность умножения. Выполняйте вычисления на бумаге.
\subsection{Как правильно поставить запятую}
\begin{enumerate}
\item Умножьте числа как целые, временно не учитывая запятые.
\item Сложите количество цифр после запятой в обоих множителях.
\item Отсчитайте столько цифр справа в произведении. Если цифр не хватает, допишите нули слева и поставьте запятую.
\end{enumerate}
\ex{Пример 3. Одинаковые цифры, разные значения}
\[
24\cdot121=2904.
\]
\begin{center}
\renewcommand{\arraystretch}{1.35}
\begin{tabular}{lcc}
\toprule
Произведение & Знаков после запятой & Результат\\
\midrule
$0,024\cdot0,121$ & $3+3=6$ & $0,002904$\\
$0,24\cdot0,121$ & $2+3=5$ & $0,02904$\\
\bottomrule
\end{tabular}
\end{center}
В обоих случаях получается одно и то же целое произведение, но положение запятой различно. Числа отличаются в десять раз.
\subsection{Деление и равные дроби}
Чтобы убрать десятичные запятые в частном, умножайте числитель и знаменатель на одно и то же ненулевое число. Записывайте оба умножения явно:
\[
\frac{11,8}{2,5}=\frac{11,8\cdot10}{2,5\cdot10}=\frac{118}{25}
=\frac{118\cdot4}{25\cdot4}=\frac{472}{100}=4,72.
\]
Если в числителе есть ещё множители, сохраняйте их:
\[
\frac{4\cdot11,8}{2,5}=\frac{4\cdot118}{25}
=\frac{4\cdot118\cdot4}{25\cdot4}=\frac{1888}{100}=18,88.
\]
\textbf{Самопроверка:} записывайте промежуточные произведения, проверяйте перенос разрядов и положение запятой. При делении убедитесь, что знаменатель не равен нулю.

\clearpage
\section{Нахождение неизвестного члена пропорции}
Примените основное свойство пропорции, а затем решите полученное уравнение.
\[
\frac{a}{b}=\frac{c}{x}\quad\Longrightarrow\quad ax=bc,
\qquad b\ne0,\ x\ne0.
\]
Если $a\ne0$, то $x=\dfrac{bc}{a}$; найденное значение должно удовлетворять исходному ограничению $x\ne0$.

Общее правило: если $A\cdot x=B$ и $A\ne0$, то $x=B:A$.
\ex{Пример 4. Пропорция со смешанными числами}
\[
\mix{3}{1}{2}:\mix{2}{1}{7}=\mix{2}{1}{3}:t,\qquad t\ne0.
\]
Перемножаем крайние и средние члены:
\[
\mix{3}{1}{2}\cdot t=\mix{2}{1}{7}\cdot\mix{2}{1}{3}.
\]
Переводим смешанные числа в неправильные дроби. Для этого умножаем целую часть на знаменатель, прибавляем числитель, знаменатель сохраняем.
\[
\mix{3}{1}{2}=\frac72,\qquad \mix{2}{1}{7}=\frac{15}{7},\qquad \mix{2}{1}{3}=\frac73.
\]
Сокращаем общие множители до умножения:
\[
\frac72\cdot t=\frac{15}{7}\cdot\frac73=\frac{15}{3}=5.
\]
Чтобы найти неизвестный множитель, делим на коэффициент при $t$:
\[
t=5:\frac72=5\cdot\frac27=\frac{10}{7}=\mix{1}{3}{7}.
\]
\textbf{Проверка:} $t\ne0$ и $\dfrac72\cdot\dfrac{10}{7}=5$ --- произведения совпали.

\textcolor{darkaccent}{Важно.} При делении на дробь переворачиваем только делитель. Сокращаем общие множители числителя и знаменателя, а не произвольные числа.

\clearpage
\section{Прямая пропорциональность}
Две положительные величины прямо пропорциональны, если при увеличении одной в несколько раз другая увеличивается во столько же раз (и аналогично при уменьшении).
\[
y=kx,\qquad \frac{y}{x}=k,\qquad k>0\text{ --- постоянное число}.
\]
\textbf{Признак в задаче:} одинаковая цена за единицу, одинаковая масса одного метра, постоянная производительность или одинаковая доля выхода продукта.

\ex{Пример 5. Семена и масло}
Из $3,5$ кг семян получают $1,4$ кг масла. Сколько масла получится из $9,5$ кг семян при том же выходе масла?
\begin{center}
\begin{tabular}{cc}
\toprule
Семена, кг & Масло, кг\\\midrule
$3,5$ & $1,4$\\
$9,5$ & $x$\\\bottomrule
\end{tabular}
\qquad
\begin{tikzpicture}[baseline=-2pt]
\draw[-{Latex},accent,thick] (0,.55)--(0,-.55);
\draw[-{Latex},accent,thick] (.65,.55)--(.65,-.55);
\end{tikzpicture}
\end{center}
Чем больше семян, тем больше масла при постоянном выходе. Стрелки сонаправлены: отношения записываем в одинаковом порядке.
\[
\frac{3,5}{9,5}=\frac{1,4}{x}
\quad\Longrightarrow\quad 3,5x=9,5\cdot1,4=13,3.
\]
\[
x=\frac{13,3}{3,5}=3,8\text{ кг}.
\]
Проверка: $3,8>1,4$, а выход масла одинаков: $1,4:3,5=3,8:9,5=0,4$.

\ex{Пример 6. Длина и масса трубы}
Однородная труба длиной $2,5$ м имеет массу $11,8$ кг. Найдём массу $x$ трубы длиной $4$ м при той же массе одного метра:
\[
\frac{2,5}{4}=\frac{11,8}{x},\qquad x=\frac{4\cdot11,8}{2,5}=18,88\text{ кг}.
\]
Значение увеличилось, как и длина трубы.

Другие примеры: масса сыра и стоимость при постоянной цене за килограмм; число кошек и расход корма за одинаковое время при одинаковой норме кормления.

\clearpage
\section{Обратная пропорциональность}
Две положительные величины обратно пропорциональны, если при увеличении одной в несколько раз другая уменьшается во столько же раз.
\[
y=\frac{k}{x},\qquad xy=k,\qquad k>0\text{ --- постоянное число}.
\]
\textbf{Признак в задаче:} одинаковое расстояние или один и тот же объём работы при неизменной производительности каждой машины.
\ex{Пример 7. Скорость и время}
При скорости $60$ км/ч автомобиль преодолевает путь за $4$ ч. Сколько времени займёт тот же путь при скорости $80$ км/ч?
\begin{center}
\begin{tabular}{cc}
\toprule
Скорость, км/ч & Время, ч\\\midrule
$60$ & $4$\\
$80$ & $x$\\\bottomrule
\end{tabular}
\qquad
\begin{tikzpicture}[baseline=-2pt]
\draw[-{Latex},accent,thick] (0,.55)--(0,-.55);
\draw[-{Latex},accent,thick] (.65,-.55)--(.65,.55);
\end{tikzpicture}
\end{center}
Расстояние неизменно. Скорость увеличивается, время уменьшается. Стрелки разнонаправлены: порядок чисел во втором отношении меняем.
\[
\frac{60}{80}=\frac{x}{4}
\quad\Longrightarrow\quad80x=60\cdot4,
\qquad x=3\text{ ч}.
\]
Проверка: $3<4$ и $60\cdot4=80\cdot3=240$ км.
\ex{Пример 8. Количество машин и время работы}
Пять одинаковых машин выполняют работу за $24$ мин. За сколько минут выполнят её восемь таких машин, работающих одновременно?
\[
\frac58=\frac{x}{24},\qquad 8x=5\cdot24,\qquad x=15\text{ мин}.
\]
Проверка: $15<24$ и $5\cdot24=8\cdot15=120$.
\subsection{Не путайте направление и пропорциональность}
Само по себе увеличение или уменьшение величин ещё не доказывает пропорциональность. Обязательно выясните, что остаётся постоянным: отношение или произведение.

\clearpage
\section{Алгоритм решения текстовой задачи}
\begin{center}
\begin{tikzpicture}[
term/.style={ellipse,draw=accent,minimum height=9mm,inner sep=4pt,align=center,font=\small},
action/.style={rectangle,rounded corners=2pt,draw=accent,text width=9cm,minimum height=12mm,align=center,font=\small,inner sep=5pt},
choice/.style={diamond,aspect=2.3,draw=accent,text width=3cm,align=center,font=\small,inner sep=2pt},
branch/.style={action,text width=4.0cm},
arr/.style={-{Latex[length=2.4mm]},draw=darkaccent,line width=.65pt}]
\node[term] (start) at (0,0) {Прочитайте условие};
\node[action] (table) at (0,-2) {Выделите две величины. Запишите их в два столбца, каждый случай --- в отдельную строку. Неизвестное обозначьте $x$. Столбцы обозначим $a$ и $b$, строки --- индексами 1 и 2.};
\node[choice] (q1) at (0,-5) {Отношение величин постоянно?};
\node[branch] (direct) at (-5.7,-5) {Прямая пропорциональность\\$\dfrac{a_1}{a_2}=\dfrac{b_1}{b_2}$\\Одинаковый порядок};
\node[choice] (q2) at (0,-8.5) {Произведение величин постоянно?};
\node[branch] (inverse) at (5.7,-8.5) {Обратная пропорциональность\\$\dfrac{a_1}{a_2}=\dfrac{b_2}{b_1}$\\Обратный порядок};
\node[term,text width=3.5cm] (other) at (0,-11.1) {Нужен другой метод};
\node[action] (solve) at (0,-13.6) {Перемножьте крайние и средние члены. Решите уравнение, учитывая ограничения.};
\node[action] (check) at (0,-16) {Проверьте вычисления, единицы измерения, постоянное отношение или произведение и смысл результата.};
\node[term] (end) at (0,-18) {Запишите ответ};
\draw[arr] (start)--(table);\draw[arr] (table)--(q1);
\draw[arr] (q1.west)--node[above,font=\small]{да}(direct.east);
\draw[arr] (q1.south)--node[right,font=\small]{нет}(q2.north);
\draw[arr] (q2.east)--node[above,font=\small]{да}(inverse.west);
\draw[arr] (q2.south)--node[right,font=\small]{нет}(other.north);
\draw[arr] (direct.south) |- (solve.west);
\draw[arr] (inverse.south) |- (solve.east);
\draw[arr] (solve)--(check);\draw[arr] (check)--(end);
\end{tikzpicture}
\end{center}

\clearpage
\section*{Тренировочный блок}
Выполняйте вычисления на бумаге. В текстовых задачах оформляйте данные в два столбца, определяйте вид зависимости и проверяйте смысл ответа.
\subsection*{Базовый уровень}
\begin{enumerate}
\item Вычислите: $2,4\cdot0,35$.
\item Вычислите: $0,018\cdot0,024$.
\item Можно ли составить пропорцию $0,036:0,12=0,9:3$?
\item Верна ли пропорция $0,036:0,12=0,9:0,3$?
\end{enumerate}
\subsection*{Средний уровень}
\begin{enumerate}[start=5]
\item Найдите неизвестный член пропорции: $1,8:2,4=3,6:x$.
\item Найдите $t$: $\mix{2}{1}{4}:\mix{1}{1}{2}=\mix{3}{1}{3}:t$.
\item Из $4,5$ кг семян получают $1,8$ кг масла. Сколько масла получится из $12,5$ кг семян при том же выходе?
\item Однородная труба длиной $2,4$ м имеет массу $10,8$ кг. Какова масса $3,7$ м такой трубы?
\end{enumerate}
\subsection*{Уровень выше среднего}
\begin{enumerate}[start=9]
\item Автомобиль проехал некоторый путь со скоростью $72$ км/ч за $5$ ч. За какое время он преодолеет тот же путь со скоростью $90$ км/ч?
\item Семь одинаковых машин, работая одновременно, выполняют заказ за $48$ мин. За какое время выполнят тот же заказ двенадцать таких машин?
\end{enumerate}

\clearpage
\section*{Ответы}
\begin{center}
\renewcommand{\arraystretch}{1.25}
\begin{tabularx}{\linewidth}{@{}rX@{}}
\toprule
№ & Ответ\\\midrule
1 & $0,84$\\
2 & $0,000432$\\
3 & Да\\
4 & Нет\\
5 & $x=4,8$\\
6 & $t=\dfrac{20}{9}=\mix{2}{2}{9}$\\
7 & $5$ кг\\
8 & $16,65$ кг\\
9 & $4$ ч\\
10 & $28$ мин\\
\bottomrule
\end{tabularx}
\end{center}
\subsection*{Итоговая памятка}
\begin{center}
\renewcommand{\arraystretch}{1.18}
\begin{tabularx}{\linewidth}{@{}p{.36\linewidth}X@{}}
\toprule
Сигнал в задаче & Полезное действие\\\midrule
Два отношения & Сравните произведения крайних и средних членов.\\
Десятичные дроби & Проверьте число знаков после запятой.\\
Смешанные числа & Переведите в неправильные дроби.\\
Неизвестный множитель & Разделите произведение на известный множитель.\\
Прямая зависимость & Сохраните порядок во втором отношении.\\
Обратная зависимость & Поменяйте порядок во втором отношении.\\
Получен ответ & Проверьте постоянное отношение или произведение и единицы.\\
\bottomrule
\end{tabularx}
\end{center}
\end{document}
